\section{任务定义与 TrajGap-Bench 评测协议}
\label{sec:bench}

\subsection{不完整历史下的直接预测任务}
\label{sec:task}

设 $t$ 时刻场景中有 $N$ 个行人。对行人 $i$，历史由位置序列 $\{p_{i,t}\}_{t=1}^{T}$ 与有效性序列 $\{m_{i,t}\}_{t=1}^{T}$ 共同给出，其中 $p_{i,t}\in\mathbb{R}^{2}$ 为位置，$m_{i,t}\in\{0,1\}$ 为有效性标记。沿用 ETH/UCY 的常规设置，取 $T=8$ 帧、帧间隔 $\Delta t=0.4$~s（历史 3.2~s），预测长度 $H=12$ 帧（4.8~s）。每个样本至少包含一个有效观测帧，不满足该条件的样本在加载阶段被移除。当 $m_{i,t}=0$ 时，位置张量对应行存放占位值，不得被解释为坐标。

直接预测器定义为
\begin{equation}
\{(\hat{y}_{i,k},\pi_{i,k})\}_{k=1}^{K}=f_{\theta}\big(\{p_{i,t},m_{i,t}\}_{i,t},\,\mathcal{C}\big),
\label{eq:task}
\end{equation}
它把带掩码的历史与场景上下文 $\mathcal{C}$（其他智能体、其掩码及其几何关系）映射为 $K$ 条候选未来 $\hat{y}_{i,k}\in\mathbb{R}^{H\times 2}$ 及模式概率 $\pi_{i,k}\in[0,1]$，且 $\sum_{k}\pi_{i,k}=1$。定义该设置的约束是流程性的而非结构性的：$f_{\theta}$ 可以执行任何内部计算（包括隐式补全），但不得调用外部插补器并把其输出替换为缺失坐标。这使直接预测与“先插补后预测”得以区分——后者的预测器接收的是稠密序列，因而严格拥有更多信息。两类设置分表报告，从不混列。

由于原始场景坐标的原点与朝向各不相同，每个样本都被表达在以目标行为中心的局部坐标系中：原点取该行人的最后一个有效观测位置，$x$ 轴取“倒数第二个有效观测位置指向最后一个有效观测位置”的方向。两个下标均在\emph{有效观测帧}上选取，因此当历史被缺失打断时，朝向估计会跨越缺口。场景内全部智能体被变换到同一坐标系，未来目标亦在其中表达。这是 ETH/UCY 的标准输入归一化方式，区别在于参考方向按缺口感知方式定义，从而避免从占位行推导朝向。

\subsection{数据来源与 Mixed 协议}
\label{sec:mixed}

TrajGap-Bench 建立在公开的 TrajImpute 数据发布之上 \cite{chib2024pedestrian}。该发布为五个 ETH/UCY 场景 \cite{lerner2007crowds,pellegrini2009you}（ETH、HOTEL、UNIV、ZARA1、ZARA2）分别提供训练、验证、测试划分，并提供由同一批底层轨迹生成的 Easy 与 Hard 两个难度版本。本文不重新生成轨迹，也不重新划分数据。对每个场景与每个划分，基准把两个版本按 actor 级样本拼接为单一的 Mixed 集合：
\begin{equation}
\text{Mixed} = \text{Easy} \cup \text{Hard},
\end{equation}
并分别作用于训练、验证与测试。该合并是样本级拼接：不重新切窗、不重采样、不去重。Easy 与 Hard 标签仅保留用于内部审计，从不对模型暴露；对外协议统一称为 Mixed。

这一选择对基准而言有三点后果，本文认为它们都是必要的。第一，训练阶段看到缺失帧数的完整范围（$0$ 至 $7$），而不是某个人为筛选的难度等级。第二，所有方法由同一个验证集决定检查点，模型无法在“自己恰好表现最好的难度”上被选中。第三，测试集保留官方的缺失程度扫描，即同一条潜在轨迹会以多个不同掩码出现。本文保留该扫描，因为删去它会悄然改变评测目标；但这意味着测试样本之间并不独立：两行样本可能描述同一条底层轨迹。因此统计比较必须按轨迹聚类，或在源轨迹层面做配对自助重采样；第~\ref{sec:disc} 节将进一步讨论。

\begin{table}[t]
\centering
\small
\caption{TrajGap-Bench 的 Mixed 划分。测试集保留官方缺失程度扫描，因此表中为样本数而非唯一源轨迹数。}
\label{tab:data}
\begin{tabular}{@{}lrrr@{}}
\toprule
Scene & Train & Val & Test \\
\midrule
ETH-M   & 59{,}618 & 10{,}698 & 1{,}629 \\
HOTEL-M & 58{,}304 & 10{,}272 & 9{,}477 \\
UNIV-M  & 18{,}462 & 5{,}416  & 219{,}006 \\
ZARA1-M & 56{,}020 & 10{,}236 & 20{,}277 \\
ZARA2-M & 51{,}014 & 8{,}346  & 52{,}497 \\
\midrule
Total   & 243{,}418 & 44{,}968 & 302{,}886 \\
\bottomrule
\end{tabular}
\end{table}

表~\ref{tab:data} 给出合并后的样本规模。样本分布的不均衡继承自原始发布：ETH-M 仅占测试样本的 0.5\%，UNIV-M 占 72\%。为此，基准同时报告五场景宏平均（各场景等权）与按样本数加权的微平均（反映测试集分布）。二者都不是综合分数，也都不替代逐场景表格。

\subsection{评测约定}
\label{sec:contract}

为使任意方法都能被纳入，评测约定被刻意保持精简。
\begin{enumerate}
\item \textbf{训练与验证。}所有方法在 Mixed-train 上训练，并仅在 Mixed-val 上选择一个检查点。监控量为最终预测头在验证集上的 $\mathrm{minFDE}_{20}$；不得依据测试数据选择轮次，也不得逐场景分别选点。
\item \textbf{测试。}Mixed-test 仅用所选检查点评估一次；不允许测试时调参、阈值化或逐场景挑选检查点。
\item \textbf{指标。}两项指标定义排序。$\mathrm{minADE}_{20}$ 为在 $K=20$ 条候选中平均位移误差最小的那条候选的平均位移误差；$\mathrm{minFDE}_{20}$ 为终点误差最小的那条候选的终点位移误差。单位均为米，越低越好。二者必须来自同一批 $K=20$ 候选、同一检查点与同一测试划分，且从不被合成为单一分数。
\item \textbf{汇报。}结果同时以逐场景值、五场景宏平均、按样本数加权的微平均与最差场景报告。未完成全部五个场景的方法不参与聚合，部分场景结果不得作为基准结论报告。
\item \textbf{设置分离。}直接预测与先插补后预测的输入不同；引入外部插补器的方法在单独设置中评估。模型内部的隐式补全属于直接预测设置的一部分，被视为模型组件。
\end{enumerate}

$K=20$ 的选择基于三点：数据发布、插补基线以及本文对比中最强的外部预测器均报告 best-of-20；在拥挤场景中，较小的 $K$ 会使 best-of-$K$ 迅速饱和；且 $K=20$ 与后续工作具有可比性。本文另报告 top-1 指标（$\mathrm{ADE@1}$ 与 $\mathrm{FDE@1}$，取概率最高的模式）作为诊断，不参与排序——输出无权重样本的方法无法计算该指标，且如第~\ref{sec:coverage} 节所示，它可以与主指标反向变化。

\subsection{分层诊断}
\label{sec:diagnostics}

每个场景仅一个数值会掩盖基准本应暴露的结构。因此本文报告五个分层维度，全部由有效性掩码与真值计算，并逐样本发布。
\begin{itemize}
\item \textbf{缺失帧数} $m_{i}=\sum_{t}(1-m_{i,t})$，历史中缺失帧的数量，取值 $[0,7]$，是粗粒度的严重程度轴。
\item \textbf{末端缺失} $\mathbb{1}[m_{i,T}=0]$，用于区分“缺口在预测起点之前结束”与“缺口一直延伸到预测起点”两类样本。
\item \textbf{锚点间隔} $T-\ell_{i}$，其中 $\ell_{i}$ 为最后一个有效观测帧的下标。当末帧被观测时为零；随最新可信状态变旧而增大。
\item \textbf{最长连续缺失}，历史中最长连续缺口的长度，用于区分零散丢帧与单次遮挡。
\item \textbf{有效帧数} $T-m_{i}$，有效观测帧的数量，作为完整性参考一并报告。
\end{itemize}

缺失帧数相同的两个样本可能在其他所有维度上都不同：六帧缺失既可能是窗口末端的一段连续缺口，也可能是散布在窗口中的六个孤立点且末帧被观测。由于速度、朝向以及最后一个状态的可信程度取决于后一种结构而非计数本身，本文把计数视为多个轴之一。所有分层均为事后分析，从不参与检查点选择或排序。

\subsection{适用范围与前提说明}
本基准是受控的压力测试，而非对真实跟踪器缺失分布的测量。掩码来自数据发布，其边缘分布不对应任何特定的检测器或跟踪器；外部效度须借助真实感知流程产出的轨迹另行建立 \cite{fujii2024realtraj}。基准亦限定于第三视角场景、五个 ETH/UCY 场景与单一历史长度。这些限制是有意为之：它们使评测契约可审计，并界定本文实证结论成立的边界。
